\scanpage{71}
Paradigme des courbes algébriques

I\quad Point de vue topologique.

On considère des surfaces topologiques (sans bord) orientables paracompactes \inserted{$X$} -- on se donne aussi un ensemble \inserted{$\omega$} de cardinal $2$, et on pose
\[
T=\mathbb Z\mathbin{\wedge}_{\mathbb Z/2}\omega\ ,
\]
et on suppose donnée pour $X$ une $\omega$-orientation, i.e. une \inserted{$\mathbb Z/2$-}application \deleted{\unsure{compatible}}\inserted{de $\omega$ dans} l'ens des orientations de $X$, ou ce qui revient au $=$, un isomorphisme du faisceau des entiers tordus sur $X$ avec le faisceau constant $T_X$. On suppose que \deleted{\unsure{les}}tte composante\deleted{s} connexe\deleted{s} $X_i$ de $X$ \deleted{sont}\inserted{est} homéomorphes à un $\hat X_i\setminus S_i$, où $\hat X_i$ est une surface compacte, et $S_i$ une partie finie de $\hat X_i$. On sait que $(\hat X_i,S_i)$ est déterminé à isom unique près par $X_i$ (comme le \unsure{décomp.|découpe|découpé} du compactifié de Čech
% archive sheet number: 70
